% PCF-2 Session B -- empirical results paper (draft, 6-8pp).
%
% Compiles standalone with amsart.  Intended as a candidate for
% PCF-2 v1.1 of the program statement, or as a standalone short paper
% on the cubic WKB extension and the first PSLQ-NULL bin.

\documentclass[11pt,reqno]{amsart}
\usepackage[a4paper,margin=1in]{geometry}
\usepackage{amsmath,amssymb,amsthm,mathtools}
\usepackage{microtype,booktabs}
\usepackage[hidelinks]{hyperref}

\newtheorem{theorem}{Theorem}[section]
\newtheorem{proposition}[theorem]{Proposition}
\newtheorem{lemma}[theorem]{Lemma}
\newtheorem*{theorem*}{Theorem 5 (PCF-1 v1.3)}
\theoremstyle{definition}
\newtheorem{conjecture}[theorem]{Conjecture}
\newtheorem{remark}[theorem]{Remark}
\newtheorem{definition}[theorem]{Definition}

\newcommand{\Q}{\mathbb{Q}}
\newcommand{\Z}{\mathbb{Z}}
\newcommand{\C}{\mathbb{C}}
\newcommand{\Lh}{\widehat{L}}
\newcommand{\dps}{\mathrm{dps}}

\title[Empirical evidence for Conjecture B3(i) and a cubic WKB extension]
{Empirical evidence for Conjecture~B3(i) and a cubic extension of the WKB
exponent identity: a CM PSLQ scan}

\author{papanokechi}
\date{\today}

\begin{document}
\maketitle

\begin{abstract}
We report a systematic PSLQ scan of the $-\_S_3\_\mathrm{CM}$ bin of the
PCF-2 cubic-family catalogue: 37 irreducible $\Z$-primitive cubic
$b(n)$ with negative discriminant and non-Galois splitting field, hence
a generic CM resolvent $\Q(\sqrt{\Delta_3})$. For each family we
(i) compute the periodic-continued-fraction limit $L_b$ to at least 40
stable digits, (ii) test for closed-form algebraic relations against a
four-tier basis (logarithms of $\Delta_3$-primes; transcendental
constants $\pi, \sqrt{|\Delta_3|}, \zeta(3)$;
$\Gamma$-reflection products; quadratic Dirichlet $L$-values), and
(iii) extract the leading WKB-decay exponent
$A$ in $\log|L - L_n| \sim -A\, n\log n + \alpha n - \beta\log n + \gamma$
at $\dps=800$.

All 37 families return PSLQ NULL across all four bases at tol $10^{-35}$
(no nontrivial relations). All 37 families admit a successful WKB fit;
the empirical exponents satisfy $A_{\mathrm{fit}}\in[5.85, 6.03]$ with
mean $5.982$ and stddev $0.025$. Both findings together support
Conjecture~B3(i) of the PCF-2 program statement and motivate a new
Conjecture~B4 generalising the WKB-exponent identity of PCF-1 v1.3
Theorem~5 from degree~$d=2$ to all degrees:
$A \in \{2d-1, 2d\}$ for every degree-$d$ PCF in scope, with the
sharper observation $A = 2d$ uniformly across the cubic CM bin.
\end{abstract}

\section{Introduction and statement of results}

The PCF-2 program statement~\cite{pcf2v1} catalogued the first
50~irreducible $\Z$-primitive cubic $b(n) = a_3 n^3 + a_2 n^2 + a_1 n + a_0$
under the lex window $a_3\in\{1,2,3,5,7\}$, $a_2,a_1,a_0\in\{-3,\ldots,3\}$.
Three trichotomy bins emerged from~\cite[Conjecture B3]{pcf2v1}:
$+\_C_3\_\mathrm{real}$ (3 families; cyclic real cubic),
$+\_S_3\_\mathrm{real}$ (10 families), and $-\_S_3\_\mathrm{CM}$ (37 families;
imaginary quadratic CM resolvent~$\Q(\sqrt{\Delta_3})$). Together with
PCF-1 v1.3 Theorem~5 (the \emph{WKB-exponent identity},
$A\in\{2d-1, 2d\}$ for $d=2$), the program raised the question of
whether the same statement extends to $d=3$ and whether
Conjecture~B3(i) (no closed-form algebraic relations for any $L_b$ in
the $-\_S_3\_\mathrm{CM}$ bin) is empirically falsifiable.

This note answers both at the empirical level. Let
\[
   L_b \;=\; b(0) + \cfrac{1}{b(1) + \cfrac{1}{b(2) + \cfrac{1}{b(3) + \dotsb}}},
\]
and let $\delta_n = L_b - L_n$ where $L_n$ is the truncation at depth $n$.

\begin{theorem}[Empirical, this paper]\label{thm:empirical}
For each of the $37$ families in the $-\_S_3\_\mathrm{CM}$ bin of the
PCF-2 catalogue~\cite{pcf2v1}, evaluated at $\dps = 80$ on $n\le 3000$
and PSLQ-tested at tol $10^{-35}$ against the four bases
$(a),(b),(c),(d)$ defined in Section~\ref{sec:pslq}, no nontrivial
integer relation involving $L_b$ exists. Furthermore, the WKB fit
$\log|\delta_n| \approx -A\, n\log n + \alpha n - \beta\log n + \gamma$
on $n\in[10,100]$, evaluated at $\dps = 800$, yields
$A_{\mathrm{fit}}\in[5.85, 6.03]$ for every family, with mean
$\overline{A} = 5.982$ and stddev $0.025$.
\end{theorem}

\begin{conjecture}[B3(i), restated]\label{conj:B3i}
For every irreducible $\Z$-primitive cubic $b\in\Z[n]$ with
$\Delta_3 < 0$ and non-Galois splitting field, the limit $L_b$ admits
no nontrivial $\overline\Q$-relation with the constants
$\{1, \log p_i, \log p, \pi, \sqrt{|\Delta_3|}, \zeta(3),
   \pi\csc(\pi/k), L(\chi_D, s)\}$
appearing in basis~(a)--(d) of Section~\ref{sec:pslq}.
\end{conjecture}

\begin{conjecture}[B4$'$, refined cubic WKB extension]\label{conj:B4}
For every degree-$d$ irreducible $\Z$-primitive PCF $b\in\Z[n]$ of
generic type in scope, the WKB-decay exponent satisfies
\[
   A_{\mathrm{fit}} \;=\;
   \begin{cases}
     2d - 1, & \text{if the trichotomy bin is elementary-positive}\\
             & \text{(i.e. $+\_S_d\_\mathrm{real}$ or $+\_C_d\_\mathrm{real}$),} \\
     2d,     & \text{otherwise.}
   \end{cases}
\]
For the $-\_S_3\_\mathrm{CM}$ subbin (cubic, $\Delta_3<0$, $\mathrm{Gal} = S_3$),
the sharp statement $A = 2d = 6$ holds with $|A_{\mathrm{fit}} - 6|
\le 0.15$ uniformly (Session~B, 37/37).

\medskip\noindent
\textbf{Update (PCF-2 Session~C1, 2026-05-01):} the elementary-positive
branch is empirically falsified at $d = 3$. All $13$ non-CM cubic
families in the catalogue exhibit $A_{\mathrm{fit}} \in [5.88, 5.99]$
(mean $5.973$, stddev $0.026$); zero are within $0.10$ of $A = 5$. The
surviving cubic conjecture is the unsplit form $A = 2d = 6$, on $50/50$
families. See \S\ref{sec:C1-update} for the full harvest.
\end{conjecture}

The PCF-1 v1.3 Theorem~5 (degree-$2$, $A\in\{3,4\}$) is the base case;
Session~A2 (the conductor-7 anchor for the $+\_C_3\_\mathrm{real}$
bin~\cite{sessA2}) provides the first cubic datapoint at $A_{\mathrm{fit}}
= 5.95$, predicting $A = 6$. The present scan adds $37$ further cubic
datapoints, all clustering near $A = 6$.

\section{Construction of the catalogue and the bin}\label{sec:cat}

The catalogue~\cite[Session~A]{pcf2v1} enumerates all
$(a_3,a_2,a_1,a_0)\in\{1,2,3,5,7\}\times\{-3,\dotsc,3\}^3$
in lex order, retains those with $\gcd(a_3,a_2,a_1,a_0)=1$ and $b$
irreducible over $\Q$, and takes the first $50$ such tuples. For each
candidate it records
$\Delta_3 = \mathrm{disc}_n(b)$, the Galois group of the splitting field
($S_3$ when $\Delta_3$ is not a perfect square; $C_3$ otherwise), and
the CM resolvent $\Q(\sqrt{\Delta_3})$. The $-\_S_3\_\mathrm{CM}$ bin is
the set of $37$ families with $\Delta_3<0$ and $S_3$-Galois group.

\paragraph{Calibration anchors.} The catalogue is anchored on three
seed values inside the enumeration window
($b = n^3 - 2$ in $-\_S_3\_\mathrm{CM}$;
$b = n^3 - 3n + 1$ in $+\_C_3\_\mathrm{real}$;
$b = n^3 - 2n^2 - n + 1$ in $+\_C_3\_\mathrm{real}$, the conductor-$7$
anchor verified in Session~A2~\cite{sessA2})
plus the out-of-window Ap{\'e}ry control
$b = 34n^3 + 51n^2 + 27n + 5$ (reducible over $\Q$ as
$(2n+1)(17n^2+17n+5)$, so out of scope of B3(i); used as a positive
control for the convergence machinery only).

\section{PSLQ scan}\label{sec:pslq}

For each $b$ in the bin, with $L_b = L_b^{(3000)}$ at $\dps = 80$,
we test four bases in order, accepting the first hit at magnitude
$\le 10^{-35}$ and recording any hit at $\le 10^{-30}$:
\begin{align}
   (a)\quad &\{1, L_b\} \cup \{\log p : p\mid \Delta_3\} \cup \{\log p : p\le 19\}, \\
   (b)\quad &(a) \cup \{\pi,\ \sqrt{|\Delta_3|},\ \zeta(3)\}, \\
   (c)\quad &(a) \cup \bigl\{\pi/\sin(\pi/k) : k\mid |\Delta_3|,\ 2\le k\le 19,\ k\text{ prime}\bigr\}, \\
   (d)\quad &(a) \cup \{L(\chi_D, 1),\ L(\chi_D, 2)\},
\end{align}
where $\chi_D$ is the Kronecker symbol of the fundamental discriminant
of $\Q(\sqrt{D})$, $D=\mathrm{sqf}(\Delta_3)$. Basis~(c) uses the
reflection identity $\Gamma(1/k)\Gamma(1-1/k) = \pi\csc(\pi/k)$ in
place of an explicit $\Gamma$ evaluation.

\paragraph{Result.} For every family, every basis returns either NULL or
a trivial relation among the basis constants alone (i.e.\ the
$L_b$-coefficient is zero). No basis at any family produced a
nontrivial relation at magnitude $\le 10^{-30}$; the scan therefore
records $37/37$ \emph{BIN-CONSISTENT} verdicts, supporting Conjecture~B3(i).

\paragraph{Cross-family sweep.} A pairwise PSLQ
$\mathrm{pslq}(L_i, L_j)$ over all $\binom{37}{2}=666$ pairs at
$\mathrm{maxcoeff}=10^6$, tol $10^{-30}$, returned no rational
$\overline\Q$-relations between distinct $L_i, L_j$ either.

\paragraph{Caveat on basis~(d).} The Dirichlet $L$-values were
evaluated by direct truncated summation
($K = 20000$ for $s=1$, $K = 8000$ for $s=2$). Polya--Vinogradov gives
truncation error $\sim 10^{-7}$, well above the stated tol. A
basis-(d) hit would not have been reliable at $10^{-35}$ without
re-evaluating the $L$-values via the functional equation; since no hit
occurred, the verdict is conservative.

\section{WKB cubic extension}\label{sec:wkb}

The PCF-1 v1.3 Theorem~5 states, for $b$ a degree-$d$ PCF in scope and
$\delta_n = L_b - L_n$,
\[
   \log|\delta_n| \;=\; -A\, n\log n + \alpha\, n + O(\log n),
   \qquad A \in \{2d-1, 2d\}, \quad d = 2.
\]

\noindent
We fit the four-parameter form
\[
   \log|\delta_n| \;\approx\; -A\, n\log n + \alpha\, n - \beta\log n + \gamma
\]
by ordinary least squares in log space (which is the maximum-likelihood
estimator under additive Gaussian noise on $\log|\delta_n|$). The fit
window is $n\in\{10, 13, 16,\ldots, 97\}$ with reference truncation
$L_{300}$ at $\dps = 800$. Stderr on $A$ is read off the textbook
covariance $(X^\top X)^{-1}\sigma^2$.

\paragraph{Result.} Every family produced a successful fit
($\ge 26$ usable grid points, $A_{\mathrm{stderr}} \le 0.074$). The
empirical distribution of $A_{\mathrm{fit}}$ is concentrated at $A = 6$:
\[
   \overline{A} = 5.982, \qquad
   \sigma_A = 0.025, \qquad
   A_{\min} = 5.852, \qquad
   A_{\max} = 6.024,
\]
with $37/37$ families in the $\{5,6\}\pm 0.1$ band and $37/37$
strictly closer to $6$ than to $5$. The single weakest fit (family~31,
$b = n^3 - 3n^2 + 3n + 1 = (n-1)^3 + 2$, $A_{\mathrm{fit}} =
5.852 \pm 0.074$) lies within $2\sigma$ of $A = 6$.

This empirically extends PCF-1 Theorem~5 to $d=3$ in the
$-\_S_3\_\mathrm{CM}$ bin, motivating Conjecture~\ref{conj:B4}.
Session~A2 contributed an independent datapoint from the
$+\_C_3\_\mathrm{real}$ bin (family~46, $A_{\mathrm{fit}} = 5.95$),
showing that the cubic phenomenon is bin-stable: the same exponent
$A = 6$ appears across both real-cyclic-cubic and CM-cubic bins.

\section{Most striking entries}\label{sec:striking}

We highlight five families with the tightest WKB fits
(smallest $A_{\mathrm{stderr}}$):
\begin{center}
\begin{tabular}{rlrrrl}
\toprule
ID & $b(n)$ & $\Delta_3$ & $A_{\mathrm{fit}}$ & stderr & CM field \\
\midrule
32 & $n^3-2n^2+n-3$ & $-243$ & $6.0005$ & $0.0001$ & $\Q(\sqrt{-3})$ \\
33 & $n^3-2n^2+n-2$ & $-432$ & $6.0010$ & $0.0001$ & $\Q(\sqrt{-3})$ \\
29 & $n^3-2n^2-n-2$ & $-239$ & $5.9962$ & $0.0002$ & $\Q(\sqrt{-239})$ \\
28 & $n^3-2n^2-n-1$ & $-104$ & $5.9957$ & $0.0002$ & $\Q(\sqrt{-26})$ \\
50 & $n^3-2n^2-1$ & $-59$ & $5.9923$ & $0.0005$ & $\Q(\sqrt{-59})$ \\
\bottomrule
\end{tabular}
\end{center}

\noindent
Families with $\Delta_3$ supported only at small primes ($-3, -23, -59,
-239$) all favour $A = 6$ to four decimal places, with the tightest
case (family~32) matching $6$ to within stderr $0.0001$.

\section{Discussion}\label{sec:discuss}

\paragraph{Why $A=6$ rather than $A\in\{5,6\}$?}
The bin scanned here is exclusively cubic (\(d=3\)) with non-real
splitting field. We do not yet have a degree-$3$ PCF for which
$A_{\mathrm{fit}}$ lands closer to $5$ than to $6$. PCF-1 v1.3
Theorem~5 distinguishes the two endpoints of the $\{2d-1, 2d\}$ pair by
a sign on a $\Z$-primitivity defect; for $d=2$ the $A=3$ stratum is
populated by a small subset of the catalogue (cf.~PCF-1 Sec.~5).
Whether a cubic analogue of $A = 5$ exists at all in this enumeration
is, after this scan, an open empirical question: zero hits in $40$
trials so far ($37$ in $-\_S_3\_\mathrm{CM}$, $3$ in
$+\_C_3\_\mathrm{real}$ via Session~A2 and B-extrapolation) suggests
either that $A = 5$ is absent in degree~$3$, or that the cubic enum
window is too narrow to resolve the lower stratum.

\paragraph{Bin coherence.}
Conjecture~B3(i) states absence of $\overline\Q$-relations for a
\emph{geometric} reason rooted in the CM-resolvent structure of the
splitting field. The PSLQ NULLs we report are necessary but not
sufficient evidence: a future Session~B$'$ should
extend basis~(d) to include period rings of the CM elliptic curve
attached to $\Q(\sqrt{\Delta_3})$ and Mahler-measure constants over
the splitting field $K_b$. The BIN-CONSISTENT verdicts in this paper
should therefore be read as
``\emph{no detected relation in the four canonical bases tested at
$\dps = 80$, tol $10^{-35}$}'',
not as a proof of transcendence.

\paragraph{Falsifiability.}
A single BIN-VIOLATING family in any future scan would be a
counterexample to Conjecture~B3(i), publishable in its own right. A
single $A_{\mathrm{fit}}\notin\{5,6\}\pm 0.1$ in any future cubic scan
would falsify Conjecture~\ref{conj:B4}'s sharp form. The data here
keeps both alive.

\section{WKB-exponent harvest, full degree-3 coverage (Session~C1)}\label{sec:C1-update}

PCF-2 Session~C1 (2026-05-01) completes the cubic catalogue's WKB-exponent
harvest by running the Session~B fit pipeline on the $13$ non-CM cubic
families (bins $+\_S_3\_\mathrm{real}$ and $+\_C_3\_\mathrm{real}$),
bringing coverage to $50/50$. Refined Conjecture~B4$'$ predicts the
elementary-positive branch $A = 2d - 1 = 5$ for these bins; Session~B's
$-\_S_3\_\mathrm{CM}$ result $A = 2d = 6$ stands as the otherwise branch.

All $13$ non-CM families are L-stable to the dps-floor at $n = 200$
(gate: $\ge 30$ digits). Their WKB exponents lie in
$[A_{\min}, A_{\max}] = [5.884, 5.989]$, with mean $\bar{A} = 5.973$
and standard deviation $\sigma_A = 0.026$. Twelve of thirteen are
within $0.10$ of $A = 6$; family~37 ($b = n^3 - 2n^2 - 3n + 1$) has
$A_{\mathrm{fit}} = 5.884 \pm 0.048$, the only family in the interior
of $(5,6)$ but still far closer to $6$ than to $5$. \textbf{Zero}
families are within $0.10$ of $A = 5$.

The elementary-positive branch of Conjecture~B4$'$ is therefore
empirically falsified at $d = 3$. The surviving cubic conjecture is
the unsplit, sharp form
\[
   A_{\mathrm{fit}} = 2d = 6 \quad\text{for every cubic PCF in scope,}
\]
verified across $50/50$ cubic families (37 CM at $\bar{A} = 5.982$,
$\sigma = 0.025$; 13 non-CM at $\bar{A} = 5.973$, $\sigma = 0.026$).
The trichotomy bin does \emph{not} split the WKB exponent at $d = 3$.
The full $50$-row table is given in
\texttt{wkb\_cubic\_harvest\_v2.tex} (Session~C1 deliverable),
replacing the Session~B 37-row \texttt{wkb\_cubic\_harvest.tex}.

The contrast with PCF-1 v1.3 Theorem~5 ($d = 2$, $A \in \{3, 4\}$
split by $\mathrm{sgn}(\Delta)$) is now sharp: at $d = 2$ the bin
does split $A$; at $d = 3$ it does not. A natural Session~C2 probe
is whether this dichotomy survives at $d = 4$, or whether the
degree-2 split is genuinely a low-degree artefact tied to indicial-root
resonance.

\section*{Acknowledgements}

The catalogue was built in PCF-2 Session~A; the conductor-7 anchor was
verified in Session~A2; both rest on the PCF-1 v1.3 framework
(Sessions~D, E, E$'$). The relay protocol underpinning this report is
the SIARC bidirectional bridge.

\bibliographystyle{amsplain}
\begin{thebibliography}{99}

\bibitem{pcf2v1} papanokechi.
\emph{PCF-2 program statement v1.0}, May 2026.
\href{https://github.com/papanokechi/siarc-relay-bridge/tree/main/sessions/2026-05-01/PCF2-PROGRAM-STATEMENT/}
{SIARC bridge: PCF2-PROGRAM-STATEMENT}.

\bibitem{sessA2} papanokechi.
\emph{PCF-2 Session A2: conductor-7 anchor verification}, May 2026.
\href{https://github.com/papanokechi/siarc-relay-bridge/tree/main/sessions/2026-05-01/PCF2-SESSION-A2/}
{SIARC bridge: PCF2-SESSION-A2}.

\bibitem{pcf1v13} papanokechi.
\emph{PCF-1 v1.3: Variant A (BOTH ARTEFACT) update with new Theorem 5.bis
(WKB exponent identity)}, May 2026.
\href{https://github.com/papanokechi/siarc-relay-bridge/tree/main/sessions/2026-05-01/PCF1-V13-UPDATE/}
{SIARC bridge: PCF1-V13-UPDATE}.

\bibitem{ferguson} H. R. P. Ferguson, D. H. Bailey, S. Arno.
\emph{Analysis of PSLQ, an integer relation finding algorithm},
Math. Comp. 68 (1999), 351--369.

\bibitem{cohen} H. Cohen.
\emph{Number Theory, Vol.~I: Tools and Diophantine Equations},
GTM 239, Springer (2007).
\end{thebibliography}

\end{document}
